\documentclass[theme=teal]{khmermath}
\usepackage{needspace}
\usepackage{colortbl}
\usepackage{tikz}
\usetikzlibrary{calc,shapes,arrows,intersections,through,backgrounds,shadows,decorations.pathmorphing,arrows.meta}
\usepackage{fancyhdr}
\usepackage{multicol}
\usepackage{tabularx}

\geometry{left=0.85cm, right=0.85cm, top=0.75cm, bottom=1.45cm, footskip=0.65cm}

% -------------------------------------------------------------------------
% Korean Math Textbook Style Color Palette & Macros
% -------------------------------------------------------------------------
\definecolor{koreaTeal}{RGB}{0, 150, 155}       % Primary Header Teal/Cyan
\definecolor{koreaDarkTeal}{RGB}{0, 105, 110}   % Dark Teal Accent
\definecolor{koreaBeige}{RGB}{247, 244, 237}     % Concept Sidebar Beige (#F7F4EE)
\definecolor{koreaBeigeBorder}{RGB}{225, 218, 205}
\definecolor{koreaOrange}{RGB}{235, 115, 30}     % Problem Number Orange
\definecolor{koreaBlue}{RGB}{35, 110, 190}      % Standard Blue
\definecolor{koreaPurple}{RGB}{125, 50, 160}    % Purple Accent
\definecolor{koreaNumber}{RGB}{115, 125, 135}   % Problem Number Gray
\definecolor{koreaYellow}{RGB}{255, 246, 175}   % Yellow Highlight Tint

\definecolor{brandNavy}{RGB}{0, 105, 110}
\definecolor{brandBlue}{RGB}{0, 150, 155}
\definecolor{brandLightBlue}{RGB}{247, 244, 237}
\definecolor{brandCrimson}{RGB}{195, 35, 60}
\definecolor{brandCoral}{RGB}{235, 80, 75}
\definecolor{brandGold}{RGB}{215, 130, 10}
\definecolor{brandGoldBg}{RGB}{255, 252, 242}
\definecolor{brandTeal}{RGB}{0, 150, 155}
\definecolor{brandTealBg}{RGB}{244, 252, 250}
\definecolor{subtleGray}{RGB}{246, 251, 249}
\definecolor{cardBorder}{RGB}{220, 212, 198}
\definecolor{trapRed}{RGB}{210, 40, 40}
\definecolor{trapBg}{RGB}{255, 246, 246}
\definecolor{trapBorder}{RGB}{245, 200, 200}

% Custom Header & Footer setup (Fancyhdr)
\pagestyle{fancy}
\fancyhf{}
\renewcommand{\headrulewidth}{0pt}
\renewcommand{\footrule}{\color{koreaTeal!60}\hrule height 0.6pt\vspace{2.5pt}}

\fancyfoot[L]{%
  \fontsize{8.5}{8.5}\selectfont\kml \color{koreaNumber}\textbf{\thepage}\ \vline\ \color{koreaDarkTeal}\textbf{៤. សមីការដឺក្រេទី១ មានមួយអញ្ញាត}%
}
\fancyfoot[C]{%
  \fontsize{8.5}{8.5}\selectfont\kby \color{koreaDarkTeal}\textbf{\ding{72}\ បង្រៀនដោយ ៖ \textcolor{koreaDarkTeal}{ក្រូត រស្មី}\ \ding{72}}%
}
\fancyfoot[R]{%
  \fontsize{8.5}{8.5}\selectfont\kby \color{koreaDarkTeal}\textbf{ Telegram: \textcolor{brandCrimson}{070897540}}\ \vline\ \color{koreaNumber}\textbf{\thepage}%
}

% Korean Header Banner
\newcommand{\koreanheader}[2]{%
  \noindent%
  \begin{tikzpicture}[x=1cm, y=1cm]
    \fill[koreaTeal, rounded corners=12pt] (0, 0.08) -- (0, 0.92) -- (3.8, 0.92) [rounded corners=0pt] -- (4.4, 0.08) -- cycle;
    \node[white, font=\fontsize{11}{11}\selectfont\kml\bfseries] at (1.9, 0.50) {#1};
    \node[anchor=west, font=\fontsize{17}{17}\selectfont\kml\bfseries, color=koreaDarkTeal] at (4.5, 0.42) {#2};
    \node[anchor=east, font=\fontsize{8.5}{8.5}\selectfont\kby, color=gray!80] at (19.2, 0.47) {$>$ ក្លឹបសិក្សាគណិតវិទ្យា PCM};
  \end{tikzpicture}%
  \vspace{2pt}%
}

% Yellow Formula Highlight Box
\newcommand{\koreahl}[1]{\colorbox{koreaYellow}{\fontsize{8.5}{8.5}\selectfont\textbf{#1}}}

% Section Header Macro
\newcommand{\khsection}[2]{%
  \vspace{2pt}%
  \noindent%
  \begin{tikzpicture}
    \fill[rounded corners=4pt, brandNavy!7] (0,0) rectangle (19.3, 0.65);
    \fill[rounded corners=3.5pt, brandNavy] (0.05, 0.05) rectangle (0.58, 0.60);
    \node[white, font=\fontsize{9.5}{9.5}\selectfont\bfseries] at (0.31, 0.33) {#1};
    \fill[brandGold, rounded corners=1pt] (0.68, 0.13) rectangle (0.72, 0.52);
    \node[anchor=west, font=\fontsize{11}{11}\selectfont\kml, color=brandCrimson] at (0.80, 0.33) {\textbf{#2}};
  \end{tikzpicture}%
  \vspace{2pt}%
}

% Page Exercise Header
\newcommand{\khexheader}[2]{%
  \noindent%
  \begin{tikzpicture}
    \fill[rounded corners=4pt, koreaTeal!12] (0,0) rectangle (19.3, 0.58);
    \fill[rounded corners=3pt, koreaTeal] (0.05, 0.05) rectangle (0.52, 0.53);
    \node[white, font=\fontsize{9}{9}\selectfont\bfseries] at (0.28, 0.29) {\ding{45}};
    \node[anchor=west, font=\fontsize{10.5}{10.5}\selectfont\kml\bfseries, color=koreaDarkTeal] at (0.65, 0.29) {កម្រងលំហាត់អនុវត្តន៍ ៖ ផ្នែកទី #1 (#2)};
  \end{tikzpicture}%
  \vspace{2pt}%
}

% Feynman Intuition Card Box
\newtcolorbox{khboxfeynman}[1][]{
  enhanced,
  colframe=koreaOrange!85,
  colback=brandGoldBg,
  boxrule=0.7pt,
  arc=5.0pt,
  borderline west={3.5pt}{0pt}{koreaOrange},
  drop fuzzy shadow=koreaOrange!12,
  left=8pt, right=8pt, top=5pt, bottom=5pt,
  title={\fontsize{9}{9}\selectfont\kml\color{white}\ding{96}\ គំនិតគ្រឹះយល់ដឹងស៊ីជម្រៅ (Mental Model)},
  coltitle=white,
  colbacktitle=koreaOrange,
  boxed title style={
    boxrule=0pt,
    arc=3.5pt,
    left=6pt, right=6pt, top=2pt, bottom=2pt
  },
  attach boxed title to top left={yshift*=-\tcboxedtitleheight/2, xshift=10pt},
  #1
}

% General Rule Box (វិធាន)
\newtcolorbox{khboxorange}[1][]{
  enhanced,
  colframe=brandNavy!40,
  colback=brandLightBlue,
  boxrule=0.7pt,
  arc=5.0pt,
  borderline west={3.5pt}{0pt}{brandNavy},
  drop fuzzy shadow=brandNavy!15,
  left=8pt, right=8pt, top=5pt, bottom=5pt,
  title={\fontsize{9}{9}\selectfont\kml\color{white}\ding{72}\ វិធាន},
  coltitle=white,
  colbacktitle=brandNavy,
  boxed title style={
    boxrule=0pt,
    arc=3.5pt,
    left=6pt, right=6pt, top=2pt, bottom=2pt
  },
  attach boxed title to top left={yshift*=-\tcboxedtitleheight/2, xshift=10pt},
  #1
}

% Model Worked Example Box with Arrow Annotations
\newtcolorbox{khboxmodel}[1][]{
  enhanced,
  colframe=koreaTeal!60,
  colback=brandTealBg,
  boxrule=0.7pt,
  arc=5.0pt,
  borderline west={3.5pt}{0pt}{koreaTeal},
  drop fuzzy shadow=koreaTeal!12,
  left=8pt, right=8pt, top=5pt, bottom=5pt,
  title={\fontsize{9}{9}\selectfont\kml\color{white}\ding{45}\ ឧទាហរណ៍គំរូ និងការវិភាគជាជំហានៗ (អមគំនួសព្រួញពន្យល់)},
  coltitle=white,
  colbacktitle=koreaTeal,
  boxed title style={
    boxrule=0pt,
    arc=3.5pt,
    left=6pt, right=6pt, top=2pt, bottom=2pt
  },
  attach boxed title to top left={yshift*=-\tcboxedtitleheight/2, xshift=10pt},
  #1
}

% Watch Out / Trap Box
\newtcolorbox{khboxtrap}[1][]{
  enhanced,
  colframe=trapBorder,
  colback=trapBg,
  boxrule=0.7pt,
  arc=4.5pt,
  borderline west={3.5pt}{0pt}{trapRed},
  drop fuzzy shadow=trapRed!10,
  left=8pt, right=8pt, top=5pt, bottom=5pt,
  title={\fontsize{8.5}{8.5}\selectfont\kml\color{white}\ding{55}\ ចំណុចប្រុងប្រយ័ត្ន ៖ កំហុសដែលសិស្សតែងតែភ្លាត់ (Watch Out!)},
  coltitle=white,
  colbacktitle=trapRed,
  boxed title style={
    boxrule=0pt,
    arc=3pt,
    left=5pt, right=5pt, top=1.8pt, bottom=1.8pt
  },
  attach boxed title to top left={yshift*=-\tcboxedtitleheight/2, xshift=10pt},
  #1
}

% Practice Card Box
\newtcolorbox{khboxpractice}[1][]{
  enhanced,
  colframe=brandTeal,
  colback=brandTealBg,
  boxrule=0.8pt,
  arc=4.5pt,
  borderline west={3.5pt}{0pt}{brandCrimson},
  drop fuzzy shadow=brandNavy!12,
  left=8pt, right=8pt, top=5pt, bottom=5pt,
  #1
}

% Self-Check Box Macro (Fills the exercise page elegantly with genuine utility)
\newcommand{\selfcheckcard}[1]{%
  \vspace{3pt}%
  \begin{tcolorbox}[
    enhanced,
    colframe=brandNavy!30,
    colback=brandLightBlue!60,
    boxrule=0.6pt,
    arc=4pt,
    left=8pt, right=8pt, top=4pt, bottom=4pt
  ]
    \fontsize{8.5}{8.5}\selectfont\kml%
    \textbf{\color{brandNavy}\ding{51}\ ការវាយតម្លៃការយល់ដឹងផ្ទាល់ខ្លួន (#1) ៖}\quad
    $\square$\ \textbf{យល់ច្បាស់}\qquad
    $\square$\ \textbf{យល់មធ្យម}\qquad
    $\square$\ \textbf{ត្រូវអនុវត្តបន្ថែម}\hfill
    \textbf{\color{brandCrimson}ពិន្ទុ ៖} \rule{1.6cm}{0.4pt} / ១០\\[3pt]
    \textbf{\color{gray!80}ចំណាំសំខាន់ ៖} \rule{16.5cm}{0.4pt}
  \end{tcolorbox}%
}

% Korean Number Macros (Two-tone Korean Textbook Style: 01, 02, 03...)
\def\splitnum#1#2{%
  \color{koreaNumber!65}#1\color{koreaOrange}#2%
}
\newcommand{\koreannum}[2][16]{%
  \mbox{%
    \fontsize{#1}{#1}\selectfont\bfseries%
    \ifnum#2<10\relax%
      \color{koreaNumber!65}0\color{koreaOrange}#2%
    \else%
      \expandafter\splitnum#2%
    \fi%
  }\,\ %
}

% Korean Exercise Header (Big 01, 02 with Title)
\newcommand{\koreanexhead}[2]{%
  \vspace{2.5pt}%
  \noindent\koreannum[16]{#1}{\fontsize{10}{10}\selectfont\kml\bfseries\color{brandNavy} #2}\par\vspace{2.5pt}%
}

% Practice Pill Badge & Korean Header
\newcommand{\phead}[1]{%
  \tikz[baseline=-0.6ex]{%
    \node[fill=brandCrimson, rounded corners=2.5pt, inner sep=2.5pt, text=white, font=\fontsize{8}{8}\selectfont\kml\bfseries] {#1};%
  }\,\ %
}
\newcommand{\koreanphead}[2]{%
  \vspace{2.5pt}%
  \noindent\koreannum[16]{#1}\phead{ប្រតិបត្តិ}{\fontsize{10}{10}\selectfont\kml\bfseries\color{brandNavy} #2}\par\vspace{2.5pt}%
}

% Khmer Sub-question Badge (Crimson Rounded Square)
\newcommand{\khlabel}[1]{%
  \tikz[baseline=-0.6ex]{%
    \node[fill=brandCrimson, rounded corners=3.2pt, inner sep=2.2pt, minimum width=1.35em, minimum height=1.32em, text=white, font=\kml\fontsize{8}{8}\selectfont, align=center] (lbl) {\raisebox{0.2pt}{#1}};%
  }\,\ %
}

% Solution Header & Dotted Line
\newcommand{\solutionsection}{%
  \vspace{3pt}%
  \noindent{\fontsize{9}{9}\selectfont\kml\bfseries\color{brandCrimson}\ding{45}\ ដំណោះស្រាយ ៖}\par\vspace{1.5pt}%
  \noindent\tikz{\draw[thick, loosely dotted, color=gray!70] (0,0) -- (19.3,0);}%
  \vspace{3pt}%
}

\setlength{\parindent}{0pt}
\setlength{\parskip}{2pt}
\setlength{\columnsep}{18pt}
\setlength{\columnseprule}{0.5pt}
\def\columnseprulecolor{\color{koreaTeal!40}}

\begin{document}

% =========================================================================
% SPREAD 1: ផ្នែកទី ១ - សមីការដឺក្រេទី១មានមួយអញ្ញាត
% PAGE 1: គំនិតគ្រឹះបែប Feynman + វិធាន + ឧទាហរណ៍គំរូ + ចំណុចប្រុងប្រយ័ត្ន
% =========================================================================
\koreanheader{មេរៀនទី ៤}{សមីការដឺក្រេទី១ មានមួយអញ្ញាត}

\khsection{១}{សមីការដឺក្រេទី១មានមួយអញ្ញាត}

\begin{khboxfeynman}
  \begin{minipage}[c]{0.68\textwidth}
    \textbf{គំនិតជញ្ជីងថ្លឹង (The Balance Scale Model) ៖}
    \begin{itemize}[leftmargin=12pt, itemsep=2pt, topsep=1pt, parsep=0pt]
      \item សមីការប្រៀបដូចជា \textbf{ជញ្ជីងដែលមានតុល្យភាព} (អង្គឆ្វេង = អង្គស្ដាំ)។
      \item $x$ គឺជា \textbf{ប្រអប់អាថ៌កំបាំង} ដែលផ្ទុកទម្ងន់ដែលយើងមិនទាន់ស្គាល់។
      \item \textbf{គោលដៅ ៖} ធ្វើឱ្យសល់តែ $x$ មួយគត់នៅអង្គម្ខាង ដើម្បីដឹងតម្លៃរបស់វា!
      \item \textbf{វិធានមាស ៖} ដើម្បីកុំឱ្យជញ្ជីងលម្អៀង អ្វីដែលអ្នកធ្វើចំពោះអង្គម្ខាង (បូក, ដក, គុណ, ចែក) អ្នកត្រូវតែធ្វើដូចគ្នាពិតប្រាកដចំពោះអង្គម្ខាងទៀត។
    \end{itemize}
  \end{minipage}%
  \hfill
  \begin{minipage}[c]{0.29\textwidth}
    \begin{center}
      \begin{tikzpicture}[scale=0.72]
        % Pivot base
        \draw[fill=brandNavy!80] (-0.35,0) -- (0.35,0) -- (0,0.7) -- cycle;
        \draw[line width=1.5pt, brandNavy] (-1.8,0.7) -- (1.8,0.7);
        % Left pan
        \draw[line width=0.8pt, gray!80] (-1.5,0.7) -- (-1.8,0.1) (-1.5,0.7) -- (-1.2,0.1);
        \draw[fill=koreaBeigeBorder, rounded corners=1pt] (-2.0,0.05) rectangle (-1.0,0.15);
        \node[draw=brandCrimson, fill=brandCrimson!10, rounded corners=2pt, inner sep=1.5pt, font=\fontsize{7}{7}\selectfont\bfseries] at (-1.5,0.4) {$x + 2$};
        % Right pan
        \draw[line width=0.8pt, gray!80] (1.5,0.7) -- (1.2,0.1) (1.5,0.7) -- (1.8,0.1);
        \draw[fill=koreaBeigeBorder, rounded corners=1pt] (1.0,0.05) rectangle (2.0,0.15);
        \node[draw=koreaTeal, fill=koreaTeal!15, rounded corners=2pt, inner sep=1.5pt, font=\fontsize{7}{7}\selectfont\bfseries] at (1.5,0.4) {$8$};
        % Equal sign
        \node[font=\fontsize{9}{9}\selectfont\bfseries, color=brandCrimson] at (0,1.05) {$=$};
        \node[font=\fontsize{7}{7}\selectfont\kml, color=gray!90] at (0,-0.3) {ជញ្ជីងមានតុល្យភាព};
      \end{tikzpicture}
    \end{center}
  \end{minipage}
\end{khboxfeynman}

\begin{khboxorange}[title={\fontsize{9}{9}\selectfont\kml\color{white}\ding{72}\ វិធានគន្លឹះ និងក្បួនបោះអង្គ (Operation Inverses)}]
  \begin{tabularx}{\linewidth}{@{}p{5.2cm}|X@{}}
    \textbf{\color{brandNavy}\ding{228}\ លក្ខណៈសមភាពគ្រឹះ ៖} & \textbf{\color{brandCrimson}\ding{228}\ ក្បួនបោះអង្គ (ប្តូរសញ្ញា \& ប្រមាណវិធី) ៖}\\[2pt]
    $\bullet$ $a = b \iff a + c = b + c$ & 
    $\bullet$ \koreahl{\textbf{បូក ($+$)} $\;\xrightarrow{\text{បោះទៅ}}\;$ \textbf{ដក ($-$)}} \quad (ឧ. $x + 2 = 8 \iff x = 8 - 2$)\\[2pt]
    $\bullet$ $a = b \iff a - c = b - c$ & 
    $\bullet$ \koreahl{\textbf{ដក ($-$)} $\;\xrightarrow{\text{បោះទៅ}}\;$ \textbf{បូក ($+$)}} \quad (ឧ. $x - 3 = 5 \iff x = 5 + 3$)\\[2pt]
    $\bullet$ $a = b \iff ac = bc$ & 
    $\bullet$ \koreahl{\textbf{គុណ ($\times$)} $\;\xrightarrow{\text{បោះទៅ}}\;$ \textbf{ចែក ($\div$)}} \quad (ឧ. $4x = 12 \iff x = \frac{12}{4}$)\\[2pt]
    $\bullet$ $a = b \iff \dfrac{a}{c} = \dfrac{b}{c}$ \ ($c \neq 0$) & 
    $\bullet$ \koreahl{\textbf{ចែក ($\div$)} $\;\xrightarrow{\text{បោះទៅ}}\;$ \textbf{គុណ ($\times$)}} \quad (ឧ. $\frac{x}{3} = 4 \iff x = 4 \times 3$)
  \end{tabularx}
\end{khboxorange}

\begin{khboxmodel}
  \begin{tikzpicture}[>=Latex, every node/.style={inner sep=1pt}]
    % Left column: Model 1
    \node[anchor=north west] at (0, 0) (title1) {\textbf{\color{koreaTeal}\ding{226}\ គំរូទី ១ ៖ $4x + 2 = 14$}};
    \node[anchor=north west] at (0.3, -0.55) (m1_s1) {\large $4x + 2 = 14$};
    \node[anchor=north west] at (0.3, -1.55) (m1_s2) {\large $4x = 12$};
    \node[anchor=north west] at (0.3, -2.55) (m1_s3) {\large $\mathbf{x = 3}$};

    % Annotations Model 1
    \draw[->, thick, koreaOrange] (2.8, -0.65) to[out=0, in=0, looseness=1.5] 
      node[right=2pt, font=\fontsize{7}{8}\selectfont, fill=koreaOrange!12, rounded corners=2pt, inner sep=2pt, draw=koreaOrange!50] {
        \textbf{\color{brandCrimson}ដក $2$ ទាំងសងខាង} ($-2$)
      } (2.8, -1.65);

    \draw[->, thick, koreaOrange] (2.2, -1.65) to[out=0, in=0, looseness=1.5] 
      node[right=2pt, font=\fontsize{7}{8}\selectfont, fill=koreaTeal!12, rounded corners=2pt, inner sep=2pt, draw=koreaTeal!50] {
        \textbf{\color{koreaDarkTeal}ចែកនឹង $4$ ទាំងសងខាង} ($\div 4$)
      } (2.2, -2.65);

    \node[anchor=north west, font=\fontsize{7.5}{8.5}\selectfont\color{koreaDarkTeal}] at (0.3, -3.15) 
      {\ding{51}\ \textbf{ផ្ទៀងផ្ទាត់ ៖} $4(3) + 2 = 12 + 2 = 14$ (ពិត)};

    % Separator line
    \draw[dashed, color=koreaTeal!35, line width=0.8pt] (8.2, 0.1) -- (8.2, -3.3);

    % Right column: Model 2
    \node[anchor=north west] at (8.6, 0) (title2) {\textbf{\color{koreaTeal}\ding{226}\ គំរូទី ២ ៖ ដោះស្រាយសមីការ}};
    
    % Equation 1 with distribution arcs
    \node[anchor=north west] at (8.8, -0.75) (m2_s1) {\large $3(x - 2) + 4 = x + 10$};
    \draw[->, thick, brandCrimson] (8.95, -0.68) to[out=60, in=120, looseness=2] node[above=0pt, font=\fontsize{6}{6}\selectfont\color{brandCrimson}] {$\times$} (9.40, -0.68);
    \draw[->, thick, brandCrimson] (8.95, -0.68) to[out=75, in=105, looseness=1.5] (9.90, -0.68);

    \node[anchor=north west] at (8.8, -1.55) (m2_s2) {\large $3x - 2 = x + 10$};
    \node[anchor=north west] at (8.8, -2.35) (m2_s3) {\large $2x = 12$};
    \node[anchor=north west] at (8.8, -3.10) (m2_s4) {\large $\mathbf{x = 6}$};

    \draw[->, thick, koreaOrange] (14.2, -0.75) to[out=0, in=0, looseness=1.3] 
      node[right=2pt, font=\fontsize{6.5}{7.5}\selectfont, fill=koreaOrange!12, rounded corners=2pt, inner sep=1.5pt, draw=koreaOrange!50] {
        \textbf{\color{brandCrimson}ពន្លាត} ($3x - 6 + 4$)
      } (14.2, -1.65);

    \draw[->, thick, koreaOrange] (13.5, -1.65) to[out=0, in=0, looseness=1.3] 
      node[right=2pt, font=\fontsize{6.5}{7.5}\selectfont, fill=koreaTeal!12, rounded corners=2pt, inner sep=1.5pt, draw=koreaTeal!50] {
        \textbf{\color{koreaDarkTeal}ផ្ទេរតួ} ($3x - x = 10 + 2$)
      } (13.5, -2.45);

    \draw[->, thick, koreaOrange] (11.5, -2.45) to[out=0, in=0, looseness=1.3] 
      node[right=2pt, font=\fontsize{6.5}{7.5}\selectfont, fill=koreaOrange!12, rounded corners=2pt, inner sep=1.5pt, draw=koreaOrange!50] {
        \textbf{\color{brandCrimson}ចែកនឹង $2$} ($\div 2$)
      } (11.5, -3.20);

    \node[anchor=north west, font=\fontsize{7.5}{8.5}\selectfont\color{koreaDarkTeal}] at (13.2, -3.20) 
      {(\ding{51}\ ផ្ទៀងផ្ទាត់ ៖ ពិត)};
  \end{tikzpicture}
\end{khboxmodel}

\begin{khboxtrap}
  \begin{itemize}[leftmargin=14pt, itemsep=1.5pt, topsep=1pt, parsep=0pt]
    \item \textbf{ភ្លាត់សញ្ញាពេលផ្លាស់ប្ដូរអង្គ ៖} កាលណាប្ដូរតួពីអង្គម្ខាងទៅអង្គម្ខាងទៀត ត្រូវប្ដូរសញ្ញាជានិច្ច ($+$ ទៅ $-$, $-$ ទៅ $+$)។
    \item \textbf{សញ្ញាដកនៅពីមុខវង់ក្រចក ៖} $-(x - 4) = -x + 4$ (សញ្ញាក្នុងវង់ក្រចកត្រូវតែត្រឡប់ពីដកទៅបូក មិនមែន $-x - 4$ ទេ!)។
  \end{itemize}
\end{khboxtrap}

% =========================================================================
% PAGE 2: កម្រងលំហាត់ផ្នែកទី ១ ទាំងមូល (រក្សាទុកលំហាត់ដដែល ១០០%)
% =========================================================================
% \clearpage
\khexheader{១}{សមីការដឺក្រេទី១មានមួយអញ្ញាត}

\koreanexhead{1}{ចូរដោះស្រាយសមីការខាងក្រោម ៖}

\begin{multicols}{3}
  \noindent
  \khlabel{ក} $4x + 2 = 8$\\[4pt]
  \khlabel{ខ} $2x + 5 = 9$\\[4pt]
  \khlabel{គ} $3x - 12 = 3$\\[4pt]
  \khlabel{ឃ} $2x + 10 = 12$\\[4pt]
  \khlabel{ង} $-3x + 1 = 4$\\[4pt]
  \khlabel{ច} $-5x - 8 = 7$
  
  \columnbreak
  \noindent
  \khlabel{ឆ} $4x + 2 = x + 5$\\[4pt]
  \khlabel{ជ} $5x + 8 = 3x - 4$\\[4pt]
  \khlabel{ឈ} $3x - 22 = 3x + 3$\\[4pt]
  \khlabel{ញ} $-2x + 14 = -2(x - 7)$\\[4pt]
  \khlabel{ដ} $-3x - 20 = 4x + 22$
  
  \columnbreak
  \noindent
  \khlabel{ឋ} $-5x - 8 + 2x = 5x - 1$\\[4pt]
  \khlabel{ឌ} $3(x - 1) + 2 = x - 12$\\[4pt]
  \khlabel{ឍ} $2(x + 3) - 10 = x + 20$\\[4pt]
  \khlabel{ណ} $4(x - 5) - 14 = 2(x - 8) + 10$\\[4pt]
  \khlabel{ត} $4(2x - 1) + 3x = -2(x + 2) + 14$
\end{multicols}

\vspace{1pt}
\koreanexhead{2}{ចូរដោះស្រាយសមីការខាងក្រោម ៖}

\begin{multicols}{2}
  \noindent
  \khlabel{ក} $4(x - 5) = 2x - 14$\\[4pt]
  \khlabel{ខ} $x + 7 = -2(x + 8)$\\[4pt]
  \khlabel{គ} $3(2y - 1) = -8(6 + y)$\\[4pt]
  \khlabel{ឃ} $-(z + 12) = 5(2z - 1)$\\[4pt]
  \khlabel{ង} $n - (8 + 4n) = 2(3n - 4)$\\[4pt]
  \khlabel{ច} $(2m + 3)^2 - (2m + 5) = 4m(m + 4) + 4$
  
  \columnbreak
  \noindent
  \khlabel{ឆ} $5p - (8 - p) = 2[-4 - (3 + 5p) - 13]$\\[4pt]
  \khlabel{ជ} $-[k - (4k + 2)] = 2 + (2k + 7)$\\[4pt]
  \khlabel{ឈ} $2 - 4\{-2x + 3[-1 + 5(1 - x)] - 1\} = 0$\\[4pt]
  \khlabel{ញ} $\sqrt{3}(x - 2) - 2x\sqrt{2} = 2(x\sqrt{3} - 3 - x\sqrt{2})$\\[4pt]
  \khlabel{ដ} $6(7x + 8) - 7(6x + 4) = 5(6 + 7x) - 35(x - 1)$\\[4pt]
  \khlabel{ឋ} $8(3 - 2p) - 16(1 - p) - 3(p - 4) = 20 - 3p$
\end{multicols}

\vspace{1pt}
\koreanphead{3}{ចូរដោះស្រាយសមីការ ៖}

\begin{khboxpractice}
  \begin{tabularx}{\linewidth}{@{}X@{}X@{}X@{}}
    \khlabel{ក} $6(x - 2) = 3(x - 8)$ &
    \khlabel{ខ} $3x + 5 = 3(x + 2)$ &
    \khlabel{គ} $6x - 4 = 2 + 6(x - 1)$
  \end{tabularx}
\end{khboxpractice}

\selfcheckcard{សមីការដឺក្រេទី១}

% =========================================================================
% SPREAD 2: ផ្នែកទី ២ - សមីការដែលមានភាគបែងជាលេខ
% PAGE 3: គំនិតគ្រឹះបែប Feynman + វិធាន + ឧទាហរណ៍គំរូ + ចំណុចប្រុងប្រយ័ត្ន
% =========================================================================
% \clearpage
\khsection{២}{សមីការដែលមានភាគបែងជាលេខ}

\begin{khboxfeynman}
  \textbf{គំនិតបំបាត់ឧបសគ្គ (Clearing the Obstacles) ៖}
  \begin{itemize}[leftmargin=14pt, itemsep=2pt, topsep=1pt, parsep=0pt]
    \item \textbf{ហេតុអ្វីសិស្សខ្លាចប្រភាគ?} ពីព្រោះភាគបែងធ្វើឱ្យការបូកដកមានភាពស្មុគស្មាញ និងងាយច្រឡំ។
    \item \textbf{យុទ្ធសាស្ត្រឆ្លាតវៃ ៖} "កម្ចាត់ភាគបែងចោលភ្លាមៗ" តាំងពីជំហានដំបូងបំផុត!
    \item \textbf{តើកម្ចាត់ដោយរបៀបណា?} ស្វែងរក \textbf{ភាគបែងរួមតូចបំផុត (LCM/ព.ប.រ)} នៃគ្រប់ភាគបែងទាំងអស់ រួចគុណអង្គទាំងពីរនៃសមីការនឹងចំនួននោះ។ ពេលសម្រួលភាគបែងរួច សមីការនឹងប្រែក្លាយជាចំនួនគត់ស្រួលគិតភ្លាម!
  \end{itemize}
\end{khboxfeynman}

\begin{khboxorange}
  \textbf{ដើម្បីដោះស្រាយសមីការដឺក្រេទី១ មានមួយអញ្ញាតដែលមានភាគបែងជាលេខ គេត្រូវ ៖}
  \begin{itemize}[leftmargin=16pt, itemsep=1.5pt, topsep=1pt, parsep=0pt]
    \item[\color{brandNavy}\ding{228}] \textbf{តម្រូវភាគបែងរួម} នៃអង្គទាំងពីររបស់សមីការ រួចលុបភាគបែងចោល ។
    \item[\color{brandNavy}\ding{228}] \textbf{លើកតួណាមានអញ្ញាតទៅអង្គម្ខាង} និងតួគ្មានអញ្ញាតទៅអង្គម្ខាងទៀត ។ តួណាប្តូរអង្គត្រូវប្តូរសញ្ញា រួចធ្វើប្រមាណវិធី ។
  \end{itemize}
\end{khboxorange}

\begin{khboxmodel}
  \begin{tikzpicture}[>=Latex, every node/.style={inner sep=1pt}]
    \node[anchor=north west] at (0, 0) (sec2_title) {\textbf{\color{koreaTeal}\ding{226}\ គំរូ ៖ ដោះស្រាយសមីការ $\dfrac{2x - 1}{3} - \dfrac{x + 2}{4} = 1$}};

    \node[anchor=north west] at (0.3, -0.55) (f_s1) {\large $\dfrac{2x - 1}{3} - \dfrac{x + 2}{4} = 1$};
    \node[anchor=north west] at (0.3, -1.55) (f_s2) {\large $4(2x - 1) - 3(x + 2) = 12$};
    \node[anchor=north west] at (0.3, -2.45) (f_s3) {\large $8x - 4 - 3x - 6 = 12 \iff 5x - 10 = 12$};
    \node[anchor=north west] at (0.3, -3.35) (f_s4) {\large $5x = 22 \implies \mathbf{x = \dfrac{22}{5}}$};

    % Arrow 1: multiply LCM
    \draw[->, very thick, koreaOrange] (6.0, -0.7) to[out=0, in=0, looseness=1.3]
      node[right=3pt, font=\fontsize{7.5}{8.5}\selectfont, fill=koreaOrange!12, rounded corners=2pt, inner sep=2pt, draw=koreaOrange!50] {
        \textbf{\color{brandCrimson}គុណគ្រប់តួទាំងអស់នឹង $12$} (កម្ចាត់ភាគបែង $\text{LCM}(3,4)=12$)
      } (6.0, -1.6);

    % Arrow 2: expand
    \draw[->, thick, koreaTeal] (7.2, -1.6) to[out=0, in=0, looseness=1.3]
      node[right=3pt, font=\fontsize{7.5}{8.5}\selectfont, fill=koreaTeal!12, rounded corners=2pt, inner sep=2pt, draw=koreaTeal!50] {
        \textbf{\color{koreaDarkTeal}ពន្លាតវង់ក្រចក} (ប្រយ័ត្នសញ្ញាដក ៖ $-3(x+2) = -3x - 6$)
      } (7.2, -2.5);

    % Arrow 3: move -10
    \draw[->, thick, koreaOrange] (12.0, -2.5) to[out=0, in=0, looseness=1.3]
      node[right=3pt, font=\fontsize{7.5}{8.5}\selectfont, fill=koreaOrange!12, rounded corners=2pt, inner sep=2pt, draw=koreaOrange!50] {
        \textbf{\color{brandCrimson}បូក $10$ ទាំងសងខាង} ($5x = 12 + 10$)
      } (12.0, -3.4);
  \end{tikzpicture}
\end{khboxmodel}

\begin{khboxtrap}
  \begin{itemize}[leftmargin=14pt, itemsep=1.5pt, topsep=1pt, parsep=0pt]
    \item \textbf{ភ្លេចគុណតួដែលគ្មានភាគបែង ៖} សិស្សភាគច្រើនគុណតែតួប្រភាគ តែភ្លេចគុណតួលេខឯករាជ្យ (ដូចជាលេខ $1$ នៅអង្គស្ដាំ ត្រូវតែគុណនឹង $12$ ក្លាយជា $12$ ជាដាច់ខាត!)។
    \item \textbf{គ្រោះថ្នាក់សញ្ញាដកនៅពីមុខប្រភាគ ៖} $-\dfrac{x+2}{4}$ ពេលគុណនឹង $12$ ត្រូវសរសេរ $-3(x + 2) = -3x - 6$ (ហាមសរសេរ $-3x + 6$)។
  \end{itemize}
\end{khboxtrap}

% =========================================================================
% PAGE 4: កម្រងលំហាត់ផ្នែកទី ២ ទាំងមូល (រក្សាទុកលំហាត់ដដែល ១០០%)
% =========================================================================
% \clearpage
\khexheader{២}{សមីការដែលមានភាគបែងជាលេខ}

\koreanexhead{1}{ចូរដោះស្រាយសមីការខាងក្រោម ៖}

\begin{multicols}{2}
  \noindent
  \khlabel{ក} $\dfrac{x}{2} + \dfrac{1}{3} = \dfrac{3x}{6}$\\[6pt]
  \khlabel{ខ} $\dfrac{3x}{5} - \dfrac{2}{3} = \dfrac{x}{15}$\\[6pt]
  \khlabel{គ} $\dfrac{x}{4} - 3 = \dfrac{3x}{2}$\\[6pt]
  \khlabel{ឃ} $\dfrac{1}{2} + \dfrac{x}{4} = \dfrac{x - 1}{6}$\\[6pt]
  \khlabel{ង} $\dfrac{x + 1}{2} - \dfrac{2}{3} = \dfrac{2x - 1}{12}$
  
  \columnbreak
  \noindent
  \khlabel{ច} $\dfrac{-x + 2}{12} + 2 = \dfrac{x + 3}{4} - 1$\\[6pt]
  \khlabel{ឆ} $2x + \dfrac{4}{3} = \dfrac{x + 3}{4} - 6$\\[6pt]
  \khlabel{ជ} $\dfrac{2(x - 3)}{15} - \dfrac{x}{3} = \dfrac{2x}{5} + 5$\\[6pt]
  \khlabel{ឈ} $\dfrac{-3(x + 2)}{20} + \dfrac{2x - 1}{4} = \dfrac{-4x + 2}{10} + \dfrac{3}{10}$\\[6pt]
  \khlabel{ញ} $3(x + 2) + \dfrac{2x + 5}{15} = \dfrac{2(x + 4)}{6} + \dfrac{x - 10}{20}$
\end{multicols}

\vspace{1pt}
\koreanexhead{2}{ចូរដោះស្រាយសមីការខាងក្រោម ៖}

\begin{multicols}{2}
  \noindent
  \khlabel{ក} $\dfrac{5x}{3} + \dfrac{9}{8} = \dfrac{7x}{6} + \dfrac{3}{16}$\\[6pt]
  \khlabel{ខ} $\dfrac{2y + 1}{3} + \dfrac{5 - y}{4} = 1 - \dfrac{3y + 2}{6}$\\[6pt]
  \khlabel{គ} $\dfrac{2(3z - 8)}{9} - \dfrac{5(z - 3)}{3} = \dfrac{3(z - 8)}{4} + \dfrac{17}{36}$\\[6pt]
  \khlabel{ឃ} $\dfrac{3t - 5}{8} + \dfrac{t + 12}{12} = 1$\\[6pt]
  \khlabel{ង} $\dfrac{x + 4}{3} - \dfrac{x - 4}{5} = 2 + \dfrac{3x - 1}{15}$\\[6pt]
  \khlabel{ច} $1 + \dfrac{2n}{3} - \dfrac{1}{2} - n = \dfrac{1 - n}{3}$\\[6pt]
  \khlabel{ឆ} $3k + \dfrac{5k - 3}{2} = 37$
  
  \columnbreak
  \noindent
  \khlabel{ជ} $2.4k + \dfrac{3(x - 1.6)}{5} = \dfrac{8.4x - 4.9}{2}$\\[6pt]
  \khlabel{ឈ} $\left(\dfrac{3}{4}k - \dfrac{2}{5}\right) + \left(\dfrac{2}{3}x + 0.6\right) - \left(\dfrac{7}{12}x - 0.3\right) = 5.8$\\[6pt]
  \khlabel{ញ} $x + \dfrac{3}{2} - \dfrac{4x - 3}{4} = 1 - \dfrac{5x - 12}{6}$\\[6pt]
  \khlabel{ដ} $8(x - 2) + 3 = 7x + \dfrac{2x - 1}{2} - 3$\\[6pt]
  \khlabel{ឋ} $\dfrac{(2y + 1)^2}{4} - \dfrac{3(y + 2)^2}{6} = \dfrac{y(y - 2)}{2}$\\[6pt]
  \khlabel{ឌ} $\dfrac{(z - 2)^2}{3} + \dfrac{z(z + 3)}{5} = \dfrac{(8z + 7)(z - 1)}{15} + \dfrac{4}{5}$
\end{multicols}

\vspace{1pt}
\solutionsection

\koreanphead{3}{ចូរដោះស្រាយសមីការ ៖}

\begin{khboxpractice}
  \begin{tabularx}{\linewidth}{@{}X@{}X@{}}
    \khlabel{ក} $\dfrac{3x}{4} + 9 = \dfrac{x}{2} + 15$ &
    \khlabel{ខ} $\dfrac{x + 7}{6} + \dfrac{2x - 8}{2} = -4$
  \end{tabularx}
\end{khboxpractice}

\selfcheckcard{សមីការមានភាគបែង}

% =========================================================================
% SPREAD 3: ផ្នែកទី ៣ - សមីការដែលអង្គទី១ជាផលគុណនៃកត្តា
% PAGE 5: គំនិតគ្រឹះបែប Feynman + វិធាន + ឧទាហរណ៍គំរូ + ចំណុចប្រុងប្រយ័ត្ន
% =========================================================================
% \clearpage
\khsection{៣}{សមីការដែលអង្គទី១ជាផលគុណនៃកត្តា}

\begin{khboxfeynman}
  \textbf{អំណាចនៃលេខសូន្យ (The Zero Superpower) ៖}
  \begin{itemize}[leftmargin=14pt, itemsep=2pt, topsep=1pt, parsep=0pt]
    \item ស្រមៃថាអ្នកមានប្រអប់ពីរ $A$ និង $B$ គុណគ្នាស្មើនឹង $0$ ($A \times B = 0$)។
    \item \textbf{សំណួរ ៖} តើអាចទៅរួចទេដែលប្រអប់ទាំងពីរគ្មានលេខសូន្យសោះ?
    \item \textbf{ចម្លើយ ៖} មិនអាចដាច់ខាត! ពីព្រោះចំនួនមិនសូន្យពីរគុណគ្នាមិនដែលចេញ $0$ ឡើយ។ ដូច្នេះ ដាច់ខាតត្រូវតែមានយ៉ាងហោចប្រអប់មួយស្មើ $0$ ជានិច្ច ៖ \koreahl{$A = 0$ ឬ $B = 0$}!
    \item \textbf{គន្លឹះមាស ៖} គោលការណ៍នេះអនុញ្ញាតឱ្យយើងបំបែកសមីការដឺក្រេទី ២ ស្មុគស្មាញ ឱ្យក្លាយជាសមីការដឺក្រេទី ១ ងាយៗពីរដាច់ដោយឡែកពីគ្នា!
  \end{itemize}
\end{khboxfeynman}

\begin{khboxorange}[title={\fontsize{9}{9}\selectfont\kml\color{white}\ding{72}\ លក្ខណៈផលគុណនៃកត្តាស្មើនឹងសូន្យ}]
  ដើម្បីឱ្យផលគុណនៃកត្តា \koreahl{$A \times B = 0$} លុះត្រាតែកត្តា \koreahl{$A = 0$} ឬ \koreahl{$B = 0$} ។\\
  \textbf{ដើម្បីដោះស្រាយសមីការដឺក្រេទី ២ ៖}
  \begin{itemize}[leftmargin=16pt, itemsep=1.5pt, topsep=1pt, parsep=0pt]
    \item[\color{brandNavy}\ding{228}] សរសេរអង្គទី១ ជា \koreahl{ផលគុណកត្តានៃដឺក្រេទី ១} (ដោយដាក់ជាកត្តារួម ឬប្រើរូបមន្តស្វ័យគុណ $a^2-b^2$)។
    \item[\color{brandNavy}\ding{228}] ឱ្យកត្តានីមួយៗស្មើនឹងសូន្យ ($0$) រួចដោះស្រាយសមីការដឺក្រេទី១ នីមួយៗ ។
  \end{itemize}
\end{khboxorange}

\begin{khboxmodel}
  \begin{tikzpicture}[>=Latex, every node/.style={inner sep=1pt}]
    % Left: Model 1
    \node[anchor=north west] at (0, 0) {\textbf{\color{koreaTeal}\ding{226}\ គំរូទី ១ ៖ $(2x - 6)(x + 4) = 0$}};
    \node[anchor=north west] at (0.3, -0.55) (q_main) {\large $(2x - 6)(x + 4) = 0$};

    % Splitting arrows
    \draw[->, very thick, koreaOrange] (1.5, -0.95) -- (0.8, -1.5) 
      node[midway, left=2pt, font=\fontsize{7}{8}\selectfont\color{koreaOrange}] {\textbf{កត្តាទី ១}};
    \draw[->, very thick, koreaTeal] (3.0, -0.95) -- (4.0, -1.5) 
      node[midway, right=2pt, font=\fontsize{7}{8}\selectfont\color{koreaTeal}] {\textbf{កត្តាទី ២}};

    % Branch 1
    \node[anchor=north west, fill=koreaOrange!8, rounded corners=2pt, draw=koreaOrange!30, inner sep=2.5pt] at (0.1, -1.6) {
      \begin{tabular}{@{}l@{}}
        $2x - 6 = 0$\\
        $\xrightarrow{+6} 2x = 6 \xrightarrow{\div 2} \mathbf{x = 3}$
      \end{tabular}
    };

    % Branch 2
    \node[anchor=north west, fill=koreaTeal!8, rounded corners=2pt, draw=koreaTeal!30, inner sep=2.5pt] at (3.6, -1.6) {
      \begin{tabular}{@{}l@{}}
        $x + 4 = 0$\\
        $\xrightarrow{-4} \mathbf{x = -4}$
      \end{tabular}
    };

    \node[anchor=north west, font=\fontsize{8}{9}\selectfont\color{brandNavy}] at (0.3, -3.05) 
      {\textbf{សន្និដ្ឋាន ៖} សមីការមានឫសពីរគឺ $\mathbf{x = 3}$ ឬ $\mathbf{x = -4}$};

    % Separator line
    \draw[dashed, color=koreaTeal!35, line width=0.8pt] (7.5, 0.1) -- (7.5, -3.3);

    % Right: Model 2
    \node[anchor=north west] at (8.0, 0) {\textbf{\color{koreaTeal}\ding{226}\ គំរូទី ២ ៖ $x^2 - 9 = 0$ \ (ឬ $x^2 - 5x = 0$)}};
    \node[anchor=north west] at (8.3, -0.55) (q2_s1) {\large $x^2 - 3^2 = 0$};
    
    \draw[->, thick, brandCrimson] (10.5, -0.65) to[out=0, in=0, looseness=1.4]
      node[right=3pt, font=\fontsize{7}{8}\selectfont, fill=brandCrimson!10, rounded corners=2pt, inner sep=2pt, draw=brandCrimson!40] {
        \textbf{\color{brandCrimson}ប្រើរូបមន្ត} $a^2 - b^2 = (a-b)(a+b)$
      } (10.5, -1.55);

    \node[anchor=north west] at (8.3, -1.55) (q2_s2) {\large $(x - 3)(x + 3) = 0$};
    
    % Splitting arrows for Model 2
    \draw[->, thick, koreaOrange] (9.2, -1.95) -- (8.6, -2.45);
    \draw[->, thick, koreaTeal] (10.5, -1.95) -- (11.2, -2.45);

    \node[anchor=north west, font=\fontsize{8}{9}\selectfont] at (8.0, -2.55) {$x - 3 = 0 \implies \mathbf{x = 3}$};
    \node[anchor=north west, font=\fontsize{8}{9}\selectfont] at (11.0, -2.55) {$x + 3 = 0 \implies \mathbf{x = -3}$};

    \node[anchor=north west, font=\fontsize{7.5}{8.5}\selectfont\color{koreaDarkTeal}] at (8.3, -3.1) 
      {\ding{51}\ \textbf{ករណីកត្តារួម ៖} $x^2 - 5x = 0 \iff x(x - 5) = 0 \implies \mathbf{x=0}$ ឬ $\mathbf{x=5}$};
  \end{tikzpicture}
\end{khboxmodel}

\begin{khboxtrap}
  \begin{itemize}[leftmargin=14pt, itemsep=1.5pt, topsep=1pt, parsep=0pt]
    \item \textbf{លក្ខណៈនេះពិតតែពេលអង្គស្ដាំស្មើ ០ ប៉ុណ្ណោះ ៖} បើ $A \times B = 6$ មិនអាចទាញ $A=6$ ឬ $B=6$ បានទេ! ត្រូវតែផ្ទេរមកឆ្វេងឱ្យអង្គស្ដាំស្មើ $0$ សិន។
    \item \textbf{ហាមចែកសម្រួលអញ្ញាត $x$ ចោល ៖} បើ $x^2 = 4x$ មិនត្រូវសម្រួល $x$ ចោលចេញជា $x = 4$ ទេ ព្រោះនឹងបាត់ឫស $x = 0$! ត្រូវសរសេរ $x^2 - 4x = 0 \iff x(x - 4) = 0$។
  \end{itemize}
\end{khboxtrap}

% =========================================================================
% PAGE 6: កម្រងលំហាត់ផ្នែកទី ៣ ទាំងមូល (រក្សាទុកលំហាត់ដដែល ១០០%)
% =========================================================================
% \clearpage
\khexheader{៣}{សមីការដែលអង្គទី១ជាផលគុណនៃកត្តា}

\koreanexhead{1}{ចូរដោះស្រាយសមីការខាងក្រោម ៖}

\begin{multicols}{3}
  \noindent
  \khlabel{ក} $x(x - 3) = 0$\\[4pt]
  \khlabel{ខ} $x(x + 2) = 0$\\[4pt]
  \khlabel{គ} $2x(x - 5) = 0$
  
  \columnbreak
  \noindent
  \khlabel{ឃ} $4x(2x - 4) = 0$\\[4pt]
  \khlabel{ង} $-2x(3x - 1) = 0$\\[4pt]
  \khlabel{ច} $(x - 1)(x + 2) = 0$
  
  \columnbreak
  \noindent
  \khlabel{ឆ} $(x + 3)(2x - 3) = 0$\\[4pt]
  \khlabel{ជ} $(3x - 4)(7x - 14) = 0$\\[4pt]
  \khlabel{ឈ} $(3x + 4)(-2x - 1) = 0$
\end{multicols}

\vspace{1pt}
\solutionsection

\koreanexhead{2}{ចូរដោះស្រាយសមីការខាងក្រោម ៖}

\begin{multicols}{3}
  \noindent
  \khlabel{ក} $x^2 - 6x = 0$\\[4pt]
  \khlabel{ខ} $2x^2 - 2 = 0$\\[4pt]
  \khlabel{គ} $5x^2 + 20x = 0$\\[4pt]
  \khlabel{ឃ} $x^2 - 6x - 7 = 0$\\[4pt]
  \khlabel{ង} $x^2 + 3x - 4 = 0$\\[4pt]
  \khlabel{ច} $x^2 - 4x + 3 = 0$
  
  \columnbreak
  \noindent
  \khlabel{ឆ} $x^2 - 2x - 3 = 0$\\[4pt]
  \khlabel{ជ} $x^2 + 8x - 9 = 0$\\[4pt]
  \khlabel{ឈ} $x^2 - 7x + 12 = 0$\\[4pt]
  \khlabel{ញ} $x^2 - 2x + 3 = 0$\\[4pt]
  \khlabel{ដ} $x^2 - 12x + 32 = 0$\\[4pt]
  \khlabel{ឋ} $y^2 - 14 = 5y$
  
  \columnbreak
  \noindent
  \khlabel{ឌ} $x^2 - x = 3x + 12$\\[4pt]
  \khlabel{ឍ} $(x - 1)^2 - 4 = 0$\\[4pt]
  \khlabel{ណ} $4(x - 1)^2 - (2x - 3)^2 = 0$\\[4pt]
  \khlabel{ត} $9(2x + 3)^2 - (x - 5)^2 = 0$\\[4pt]
  \khlabel{ថ} $(3x - 1)^2 = 16(x + 2)^2$
\end{multicols}

\vspace{1pt}
\solutionsection

\koreanphead{3}{ចូរដោះស្រាយសមីការខាងក្រោម ៖}

\begin{khboxpractice}
  \begin{tabularx}{\linewidth}{@{}X@{}X@{}}
    \khlabel{ក} $x^2 + 2x - 4x + 8 = 0$ &
    \khlabel{គ} $3x^2 + 9x = 0$ \\[6pt]
    \khlabel{ខ} $(2x - 1)^2 - (x - 1)^2 = 0$ &
    \khlabel{ឃ} $2x^2 + 3x - 5 = 0$
  \end{tabularx}
\end{khboxpractice}

\selfcheckcard{ផលគុណកត្តាស្មើ ០}

% =========================================================================
% SPREAD 4: ផ្នែកទី ៤ - ចំណោទសមីការដឺក្រេទី១មានមួយអញ្ញាត
% PAGE 7: គំនិតគ្រឹះបែប Feynman + វចនានុក្រមបកប្រែ + ឧទាហរណ៍គំរូ + ចំណុចប្រុងប្រយ័ត្ន
% =========================================================================
% \clearpage
\khsection{៤}{ចំណោទសមីការដឺក្រេទី១មានមួយអញ្ញាត}

\begin{khboxfeynman}
  \textbf{វចនានុក្រមបកប្រែ «ភាសាជីវិត» ទៅជា «ភាសាគណិត» (Life to Math) ៖}
  \begin{center}
    \small
    \begin{tabularx}{0.96\linewidth}{|>{\raggedright\arraybackslash}X|>{\centering\arraybackslash}p{2.8cm}|>{\raggedright\arraybackslash}X|}
      \hline
      \rowcolor{koreaTeal!15}
      \textbf{ពាក្យក្នុងជីវិតប្រចាំថ្ងៃ} & \textbf{សញ្ញាគណិត} & \textbf{ឧទាហរណ៍បកប្រែ} \\
      \hline
      មួយចំនួនមិនស្គាល់, ចម្ងាយ, អាយុ & $x, y$ & តាង $x$ ជាចំនួនមិនស្គាល់ \\
      ច្រើនជាង, ថែម, ផលបូក, ចាស់ជាង & $+$ & ថែម $5$ លើមួយចំនួន $\to x + 5$ \\
      តិចជាង, បន្ថយ, ផលដក, ក្មេងជាង & $-$ & បន្ថយ $3$ ពីមួយចំនួន $\to x - 3$ \\
      ពីរដង, បីដង, ផលគុណ, ពាក់កណ្ដាល & $\times, \div$ & ពីរដងនៃមួយចំនួន $\to 2x$ \\
      ស្មើនឹង, ស្មើ, គឺ, សរុបជា & $=$ & ផលបូកស្មើ $20 \to x + (x+4) = 20$ \\
      \hline
    \end{tabularx}
  \end{center}
  \vspace{1pt}
  \textbf{យុទ្ធសាស្ត្រ ៤ ជំហាន ៖} \textbf{១. ស្វែងយល់} (តាង $x$) $\to$ \textbf{២. បកប្រែ} (បង្កើតសមីការ) $\to$ \textbf{៣. ដោះស្រាយ} (រក $x$) $\to$ \textbf{៤. ផ្ទៀងផ្ទាត់ និងសន្និដ្ឋាន}។
\end{khboxfeynman}

\begin{khboxmodel}
  \textbf{\color{koreaTeal}\ding{226}\ ចំណោទគំរូ ៖} ម៉ារីមានអាយុច្រើនជាងបូណា $6$ ឆ្នាំ។ ផលបូកអាយុអ្នកទាំងពីរស្មើនឹង $30$ ឆ្នាំ។ រកអាយុអ្នកទាំងពីរ។\\[3pt]
  \begin{tikzpicture}[>=Latex, every node/.style={inner sep=1pt}]
    % Step 1: Translation blocks
    \node[anchor=north west, fill=koreaTeal!10, rounded corners=2pt, draw=koreaTeal!50, inner sep=3pt] at (0, 0) (step1_a) {
      \textbf{អាយុបូណា} ៖ $x$
    };
    \node[anchor=north west, fill=koreaOrange!10, rounded corners=2pt, draw=koreaOrange!50, inner sep=3pt] at (3.5, 0) (step1_b) {
      \textbf{អាយុម៉ារី} ៖ $x + 6$ \ \fontsize{7}{8}\selectfont\color{gray}(ច្រើនជាង $6$ ឆ្នាំ)
    };

    % Arrow Step 1 to Step 2
    \draw[->, very thick, brandCrimson] (2.2, -0.55) -- (2.2, -1.25)
      node[midway, right=3pt, font=\fontsize{7.5}{8.5}\selectfont\color{brandCrimson}] {
        \textbf{បកប្រែជាសមីការ} (ផលបូកអាយុ $= 30$)
      };

    % Step 2: Equation
    \node[anchor=north west] at (0, -1.45) (eq_main) {
      \large $x + (x + 6) = 30 \iff 2x + 6 = 30$
    };

    % Arrow Step 2 to Step 3
    \draw[->, thick, koreaOrange] (7.2, -1.55) to[out=0, in=0, looseness=1.4]
      node[right=3pt, font=\fontsize{7}{8}\selectfont, fill=koreaOrange!12, rounded corners=2pt, inner sep=2pt, draw=koreaOrange!50] {
        \textbf{\color{brandCrimson}ដក $6$ សងខាង} ($2x = 30 - 6$)
      } (7.2, -2.45);

    \node[anchor=north west] at (0, -2.35) (eq_step2) {
      \large $2x = 24$
    };

    \draw[->, thick, koreaTeal] (2.2, -2.45) to[out=0, in=0, looseness=1.4]
      node[right=3pt, font=\fontsize{7}{8}\selectfont, fill=koreaTeal!12, rounded corners=2pt, inner sep=2pt, draw=koreaTeal!50] {
        \textbf{\color{koreaDarkTeal}ចែកនឹង $2$ សងខាង} ($\div 2$)
      } (2.2, -3.25);

    \node[anchor=north west] at (0, -3.15) (eq_step3) {
      \large $\mathbf{x = 12}$
    };

    % Step 4: Conclusion Box
    \node[anchor=north west, fill=koreaDarkTeal!8, rounded corners=3pt, draw=koreaDarkTeal, inner sep=4pt] at (5.5, -2.95) (conclusion) {
      \begin{tabular}{@{}l@{}}
        \textbf{\color{koreaDarkTeal}\ding{51}\ សន្និដ្ឋាន \& ឆ្លើយ ៖}\\
        $\bullet$ \textbf{អាយុរបស់បូណា ៖} $\mathbf{12\ \text{ឆ្នាំ}}$\\
        $\bullet$ \textbf{អាយុរបស់ម៉ារី ៖} $12 + 6 = \mathbf{18\ \text{ឆ្នាំ}}$\\
        \fontsize{7}{8}\selectfont\color{gray}(\ding{51}\ ផ្ទៀងផ្ទាត់ ៖ $12 + 18 = 30\ \text{ឆ្នាំ}$ ពិត)
      \end{tabular}
    };
  \end{tikzpicture}
\end{khboxmodel}

\begin{khboxtrap}
  \begin{itemize}[leftmargin=14pt, itemsep=1.5pt, topsep=1pt, parsep=0pt]
    \item \textbf{ប្រយ័ត្នរឿងខ្នាត និងឯកតា ៖} កាលណាប្រធានមានល្បឿនគិតជា $\text{km/h}$ ហើយរយៈពេលគិតជា "នាទី" ត្រូវតែប្ដូរនាទីទៅជាម៉ោងសិន (ឧ. $20$ នាទី $= \dfrac{20}{60}\text{h} = \dfrac{1}{3}\text{h}$) មុននឹងយកទៅគណនា!
    \item \textbf{កុំភ្លេចឆ្លើយសន្និដ្ឋាន ៖} រកឃើញ $x$ ហើយ ត្រូវឆ្លើយនូវអ្វីដែលប្រធានសួរ (កុំឆ្លើយត្រឹម $x=...$ ត្រូវបញ្ជាក់ថាជាអាយុ ឬជាចំនួនណា)។
  \end{itemize}
\end{khboxtrap}

% =========================================================================
% PAGE 8: កម្រងលំហាត់ផ្នែកទី ៤ ទាំងមូល (រក្សាទុកលំហាត់ដដែល ១០០%)
% =========================================================================
% \clearpage
\khexheader{៤}{ចំណោទសមីការដឺក្រេទី១មានមួយអញ្ញាត}

\koreannum[13]{1} ផលបូកនៃមួយចំនួនថែមនឹង $25$ ស្មើនឹង $42$ ។ ចូររកមួយចំនួននោះ ។\\[4pt]
\koreannum[13]{2} ពីរដងនៃមួយចំនួនថែមនឹង $15$ ស្មើនឹង $53$ ។ ចូររកមួយចំនួននោះ ។\\[4pt]
\koreannum[13]{3} បីដងនៃមួយចំនួនដក $20$ ស្មើមួយចំនួននោះថែម $12$ ។ ចូររកមួយចំនួននោះ ។\\[4pt]
\koreannum[13]{4} ផលបូកប្រាំចំនួនគត់តគ្នាស្មើនឹង $520$ ។ រកចំនួនគត់ទាំងប្រាំនោះ ។\\[4pt]
\koreannum[13]{5} កសិករមួយក្រុមកំពុងភ្ជួរស្រែប្រវាស់ដៃគ្នា ។ បើគេរាប់ទាំងគោ និងកសិករមានចំនួនសរុប $36$ តែបើគេរាប់ជើងវិញមានចំនួន $100$ ។ រកចំនួនកសិករ និងចំនួនគោ ។\\[4pt]
\koreannum[13]{6} ឧត្តមមានអាយុ $51$ ឆ្នាំ ហើយវីរៈមានអាយុ $21$ ឆ្នាំ ។ តើរយៈពេលប៉ុន្មានឆ្នាំទៀតទើបឧត្តមមានអាយុស្មើពីរដងនៃអាយុវីរៈ ?\\[4pt]
\koreannum[13]{7} បច្ចុប្បន្នកាលនេះឪពុកមានអាយុ $42$ ឆ្នាំ ហើយកូនមានអាយុ $15$ ឆ្នាំ ។ តើនៅប៉ុន្មានឆ្នាំទៀតទើបអាយុឪពុកស្មើនឹងពីរដងនៃអាយុកូន ?

\vspace{3pt}
\koreanphead{8}{ចូរដោះស្រាយចំណោទខាងក្រោម ៖}

\begin{khboxpractice}
  \noindent
  \khlabel{ក} ម៉ារីមានអាយុ $8$ ឆ្នាំច្រើនជាងអាយុបូណា ។ ផលបូកអាយុរបស់ម៉ារី និងបូណាស្មើនឹង $38$ ឆ្នាំ ។ ចូររកអាយុរបស់បូណា ។\\[4pt]
  \khlabel{ខ} ថែម $7$ លើពីរដងនៃមួយចំនួនស្មើនឹងបន្ថយ $27$ ពីបីដងនៃចំនួននោះ ។ ចូររកចំនួននោះ ។
\end{khboxpractice}

\vspace{3pt}
\noindent\tikz[baseline=-0.6ex]{%
  \node[fill=brandNavy, rounded corners=3pt, inner sep=3.2pt, text=white, font=\fontsize{9}{9}\selectfont\kml\bfseries] {លំហាត់អនុវត្ត};%
}\par\vspace{3pt}

\noindent
\koreannum[13]{9} \textbf{លំហាត់ទី១ ៖} ឧទ្ធម្ភាគចក្រមួយគ្រឿងធ្វើដំណើរចេញពីអាកាសយានដ្ឋានអន្តរជាតិភ្នំពេញឆ្ពោះទៅរួមខេត្តសៀមរាបដោយល្បឿន $288\text{km/h}$ ។ $20$ នាទីក្រោយមកមានយន្តហោះមួយគ្រឿងទៀតបានហោះចេញពីកន្លែងតែមួយទៅកាន់ទីរួមខេត្តសៀមរាបដែរក្នុងល្បឿន $480\text{km/h}$ ។ តើយន្តហោះនេះហោះអស់រយៈពេលប៉ុន្មានទើបតាមទាន់ឧទ្ធម្ភាគចក្រ ?\\[4pt]
\koreannum[13]{10} \textbf{លំហាត់ទី២ ៖} អ្នកជិះកង់ម្នាក់ចេញពីទីក្រុង $A$ ទៅទីក្រុង $B$ ដែលមានចម្ងាយ $50\text{km}$ ។ មួយម៉ោងក្រោយមកអ្នកជិះម៉ូតូម្នាក់ចេញពីទីក្រុង $A$ ដែរ ហើយទៅដល់ទីក្រុង $B$ មុនអ្នកជិះកង់មួយម៉ោង ។ ចូរគណនាល្បឿននៃយានទាំងពីរ បើគេដឹងថាល្បឿនម៉ូតូស្មើនឹង $2.5$ ដងនៃល្បឿនកង់ ។

\vspace{8pt}
\begin{center}
  \color{brandCrimson}\rule{0.22\textwidth}{0.8pt}\quad\ding{170}\ \textbf{\kml សំណាងល្អ}\ \ding{170}\quad\rule{0.22\textwidth}{0.8pt}
\end{center}

\end{document}
